\section{Analysis and Discussion}

\begin{enumerate}
    \item \textbf{Volume:}
            The dataset contains approximately 41.62GB of CSV and text data. The largest transaction datasets are HI-Large (17.05 GB) and LI-Large (16.74 GB), contain approximately 356M transactions, creating significant storage and computational requirements.
    \item \textbf{Velocity:}
            Although the data is synthetic and historical rather than real-time, HI-Large contains approximately 1.86M transactions per day, while LI-Large contains approximately 1.81M. This reflects the high-speed processing requirements of transaction-monitoring systems. 
    \item \textbf{Variety:}
            The dataset includes transaction amounts, currencies, payment methods, sending and receiving institutions, account identifiers, timestamps and pattern files describing different money-laundering schemes.
    \item \textbf{Value:}
            The dataset provides value for banks, regulators and financial crime research by supporting transaction monitoring, suspicious activity analysis and AML model development without requiring access to sensitive real-world banking data. 
    \item \textbf{Veracity:}
            The synthetic data provides known money-laundering labels, allowing controlled evaluation without exposing customer information. However, synthetic data cannot fully capture the complexity of real financial behaviour.
    \item \textbf{Variability:}
            Transaction amounts, currencies, payment formats and laundering behaviour vary across the dataset. The HI and LI dataset variants contain different proportions of laundering activity. 
    \item \textbf{Visualisation:}
            Transaction networks, time-series plots and distribution charts can reveal unusual activity, suspicious accounts clusters and differences between legitimate and laundering behaviour.
 
\end{enumerate}
 
 
 


